\scanpage{14}
\mspage{\handwritten{29}}
% One continuous oblique stroke crosses the typescript, from the head to the penultimate line.
\deleted{images dans $S$ sont des parties $S_i$ à la fois ouvertes et fermées, et elles recouvrent $S$. Si $S$ est connexe ou quasi-compact, il existe donc un ensemble fini d'indices \deleted{tel} $i$ tels que les $S_i$ recouvrent $S$, soit $S'$ la réunion des $P_i$ correspondants, \deleted{qui est donc tel que} \inserted{Alors} $S'\handwritten{\mathord{-}}S$ \deleted{\unsure{r}}est fini étale surjectif, et posant $P'=P\times_S S'=\underline{\operatorname{Isom}}_{S'\text{-gr}}(H',G')$, on voit que $P'$ a une section sur $S'$, donc il existe un isomorphisme de $S'$-groupes $H'=H\times_S S'\handwritten{\xrightarrow{\sim}}G'=G\times_S S'=D_{S'}(M)$, ce qui prouve que $S'$ splitte $H$. Le cas où \deleted{\unsure{S}}les composantes \deleted{irréductibles}\inserted{connexes} de $S$ sont \deleted{\unsure{connex}} ouvertes se ramène aussitôt au cas où $S$ est connexe. \quad - \quad Remarquons d'ailleurs qu'on peut, lorsque $H$ est de type fini sur $S$,\deleted{\unsure{connexe|connex}} et de type constant, donner \deleted{un}le critère d'isotrivialité suivant (où il n'est plus nécessaire de faire aucune restriction sur $S$): $R$ est réunion d'une \uline{suite} de parties à la fois ouvertes et fermées finies sur $S$.

\uline{Lemme} 5.1\deleted{2}\handwritten{3}. Soient $S$ un préschéma localement noethérien et connexe, $\deleted{R}\handwritten{P}$ un $S$-préschéma constant tordu quasi-isotrivial, $Z$ une partie à la fois ouverte et fermée de $\deleted{R}\handwritten{P}$, telle qu'il existe un $s\handwritten{\in}S$ tel que $\deleted{\unsure{R}} \handwritten{Z}_s$ soit fini. \margin{\handwritten{$Z_s$}} Alors $\deleted{\unsure{R}}\handwritten{Z}$ est fini sur $S$. \margin{\handwritten{$Z$}}

\deleted{Oublions un instant que $S$ est supposé connexe, et} \inserted{Ne supposons pas d'abord $S$ connexe;} soit $U$ l'ensemble des $s\handwritten{\in}S$ tels que $Z_s$ soit fini, nous allons prouver que $U$ est à la fois ouvert et fermé, et que $Z\handwritten{|}U$ est fini sur $U$. C'est là évidemment un \deleted{énoncé} énoncé essentiellement équivalent à 5.1\deleted{2}\handwritten{3}. mais qui a l'avantage d'être ``de nature locale'' sur $S$ pour la topologie étale (disons), ce qui nous ramène au cas où $\deleted{R}\handwritten{P}$ est constant, i.e. de la forme $\deleted{\unsure{K}}I_S$, où $I$ est un ensemble. (NB l'hypothèse localement noethérienne n'est pas perdue par un changement de base étale;} c'est là où on utilise la quasi-isotrivialité de $\deleted{R}\handwritten{P}$ sur $S$). On peut de plus supposer $S$ connexe\handwritten{,} \deleted{Mais alors}
% archive sheet number: 13

